\lyseis{Β2.1}{Συμμετρία ως προς άξονα}
\sol{1} Ελεύθερη κατασκευή. Η κάθετη από το $\mathrm{Α}$ στην $\varepsilon$ περνά από το μέσο του $\mathrm{ΑΑ}'$.
\sol{2} Το ίδιο το σημείο.
\sol{3} $3{,}5$ cm από την $\varepsilon$ και $\mathrm{ΑΑ}'=7$ cm.
\sol{4} Ελεύθερη κατασκευή. Βρίσκω τα συμμετρικά των τριών κορυφών και τα ενώνω.
\sol{5} Το συμμετρικό τρίγωνο έχει $\mathrm{Α}'\mathrm{Β}'=5$ cm και $\widehat{\mathrm{Β}'}=40^\circ$.
\sol{6} Το συμμετρικό τμήμα έχει μήκος $4{,}2$ cm.
\sol{7} α) Σωστό. β) Λάθος: κάθετο. γ) Σωστό. δ) Σωστό. ε) Λάθος: διατηρεί τα μήκη.
\sol{8} Στο συμμετρικό σημείο: $1{,}2$ μ. από το μονοπάτι, από την άλλη μεριά. Τα δύο δέντρα απέχουν $2{,}4$ μ.
\sol{9} Η αντανάκλαση φαίνεται $4$ μ. βαθιά, όσο ύψος έχει το δέντρο.
\sol{10} Πραγματικά είναι $9{:}00$ (ο ωροδείκτης στον καθρέφτη είναι αριστερά αντί δεξιά).
\sol{11} Μένουν ίδια τα $\mathrm{Α}$, $\mathrm{Η}$, $\mathrm{Μ}$, $\mathrm{Ο}$, $\mathrm{Τ}$.
\sol{12} Ελεύθερη κατασκευή.
\sol{13} Η δίπλωση φέρνει το $\mathrm{Α}$ πάνω στο $\mathrm{Α}'$ μόνο αν το $\mathrm{ΑΑ}'$ είναι κάθετο στον άξονα και ο άξονας το περνά από το μέσο του, αλλιώς τα δύο σημεία δεν θα έπεφταν το ένα πάνω στο άλλο.
\sol{14} Τα σημεία της ίδιας της $\varepsilon$.
\sol{15} Επιστρέφει στο $\mathrm{Α}$.
\sol{16} Το $\mathrm{Μ}$ ανήκει στην $\varepsilon$, άρα είναι συμμετρικό του εαυτού του. Άρα το $\mathrm{Α}'\mathrm{Β}'$ περνά από το $\mathrm{Μ}$.
\sol{17} Ελεύθερη απάντηση: σχηματίζεται ένα συμμετρικό σχήμα με άξονα την πτυχή.
\sol{18} Ελεύθερη απάντηση.
